% tenkz-cd — L2b: the commutative-diagram sub-language.
%
%   \tntree[keys]{(((a\,b)_x\,c)_y\,d)_e}
%       a fusion tree from a parenthesization word: leaves on a top
%       line, one junction per fused pair, internal charges on the
%       downward edges, a short rooted leg carrying the total charge.
%
%   \begin{tenkzcd}[polygon=n, radius=..]
%     cell & cell & ... \tnarrow[from=i, to=j]{label} ...
%   \end{tenkzcd}
%       polygon mode: n cells on a circle (vertex 1 at the top,
%       clockwise), annotation arrows between them.  \tnarrow* puts the
%       label on the other side of its arrow.
%
%   \begin{tenkzcd}[keys]  ... tikz-cd material ...  \end{tenkzcd}
%       grid mode (no polygon key): a thin wrapper around tikz-cd; cells
%       may hold tenkz objects (\tntree, \tnpic) since those are
%       self-contained boxes.
%
% Both environments and the standalone tree are math atoms: the baseline
% is the vertical middle of the picture lowered by the math axis, so a
% diagram centers on the axis of the surrounding formula.

\ExplSyntaxOn

% ---------- tree metrics (ratios of pitch) ----------
% Sibling pitch of the leaf row and the depth of one fusion level share
% one ratio, so every edge of a balanced pair meets its junction at 45
% degrees; the root leg is shorter than a level so the tree reads as
% rooted, not as one more leaf.
\def\tenkz@r@treestep{0.55}
\def\tenkz@r@treeleg{0.35}

% ---------- state ----------
\seq_new:N \l__tenkzcd_items_seq     % parenthesization word, one item each
\seq_new:N \l__tenkzcd_stack_seq     % parsed subtrees: {x}{level}{charge}
\tl_new:N  \l__tenkzcd_word_tl
\tl_new:N  \l__tenkzcd_item_tl
\tl_new:N  \l__tenkzcd_charge_tl
\tl_new:N  \l__tenkzcd_ca_tl
\tl_new:N  \l__tenkzcd_cb_tl
\tl_new:N  \l__tenkzcd_cha_tl
\tl_new:N  \l__tenkzcd_chb_tl
\tl_new:N  \l__tenkzcd_anchor_tl
\int_new:N \l__tenkzcd_leaf_int
\int_new:N \l__tenkzcd_la_int
\int_new:N \l__tenkzcd_lb_int
\int_new:N \l__tenkzcd_lp_int
\dim_new:N \l__tenkzcd_xa_dim
\dim_new:N \l__tenkzcd_xb_dim
\dim_new:N \l__tenkzcd_xp_dim
\dim_new:N \l__tenkzcd_off_dim
\dim_new:N \l__tenkzcd_axis_dim
% tenkzcd
\int_new:N \l__tenkzcd_ngon_int
\int_new:N \l__tenkzcd_i_int
\tl_new:N  \l__tenkzcd_radius_tl
\tl_new:N  \l__tenkzcd_pass_tl
\dim_new:N \l__tenkzcd_r_dim
\dim_new:N \l__tenkzcd_vx_dim
\dim_new:N \l__tenkzcd_vy_dim
\fp_new:N  \l__tenkzcd_angle_fp
\seq_new:N \l__tenkzcd_cells_seq
\seq_new:N \g__tenkzcd_arrows_seq    % {from}{to}{side keys}{label}
\tl_new:N  \l__tenkzcd_afrom_tl
\tl_new:N  \l__tenkzcd_ato_tl
\tl_new:N  \l__tenkzcd_body_tl

% ---------- keys ----------
\pgfkeys{
  /tenkz/cd/.is~family,
  /tenkz/cd/polygon/.code={\int_set:Nn \l__tenkzcd_ngon_int {#1}},
  /tenkz/cd/radius/.store~in=\l__tenkzcd_radius_tl,
  /tenkz/cd/pitch/.code={\pgfkeys{/tenkz/pitch=#1}},
  /tenkz/cd/compact/.code={\pgfkeys{/tenkz/compact}},
  /tenkz/cd/inline/.code={\pgfkeys{/tenkz/inline}},
  % the shared forward-to-core plumbing: the glyph policy reads the same
  % in every family, so \tnpic cells inside a cd picture obey it
  /tenkz/cd/tensor~style/.code={\pgfqkeys{/tenkz}{tensor~style=#1}},
  % unrecognized keys pass through to the underlying picture
  /tenkz/cd/.unknown/.code={
    \tl_if_eq:NNTF \pgfkeyscurrentvalue \pgfkeysnovalue@text
      { \tl_put_right:Ne \l__tenkzcd_pass_tl
          { \exp_not:o \pgfkeyscurrentname , } }
      { \tl_put_right:Ne \l__tenkzcd_pass_tl
          { \exp_not:o \pgfkeyscurrentname =
            \exp_not:o \pgfkeyscurrentvalue , } }
  },
}
\pgfkeys{
  /tenkz/arrow/.is~family,
  /tenkz/arrow/from/.store~in=\l__tenkzcd_afrom_tl,
  /tenkz/arrow/to/.store~in=\l__tenkzcd_ato_tl,
}

% ---------- the fusion tree ----------
% Grammar of the parenthesization word:
%   node   :=  leaf  |  ( node node ) [ _ charge ]
%   leaf, charge  :=  one character or one braced group
% Separators (\, and spaces) are stripped before parsing.  The word is
% flattened into a sequence of items and consumed left to right by a
% recursive-descent parser; every finished subtree pushes its root
% {x}{level}{charge} onto a value stack, so a fused pair pops its two
% children, joins them, and pushes the joint.
%
% Layout: leaves sit on the line y = 0, spaced by one tree step; an
% internal vertex sits at the x-midpoint of its children's roots, one
% level below the deeper child (level = 1 + max of the child levels,
% y = -level steps).  Charges are typeset only when the vertex's
% downward edge is known, so they can be set clear of the ink.

% y-coordinate of a fusion level (integer factor on the right)
\cs_new:Npn \tenkz_cd_levely:n #1
  { \dim_eval:n { -\tenkz@r@treestep\tenkz@pitch * (#1) } }

% ---------- picture framing ----------
% Every cd picture is a math atom sharing one baseline.  Read the axis
% height (NFSS sets math fonts up lazily, so force them first), then lower
% the picture's bounding-box centre by it so the diagram centres on the
% surrounding formula's axis.
\cs_new_protected:Npn \tenkz_cd_read_axis:
  { \check@mathfonts
    \dim_set:Nn \l__tenkzcd_axis_dim { \fontdimen22\textfont2 } }
\tikzset{
  tenkz~cd~baseline/.style={
    baseline={([yshift=-\dim_use:N \l__tenkzcd_axis_dim]
               current~bounding~box.center)} } }

\cs_new_protected:Npn \tenkz_cd_tree_node:
  {
    \seq_pop_left:NNTF \l__tenkzcd_items_seq \l__tenkzcd_item_tl
      {
        \str_if_eq:VnTF \l__tenkzcd_item_tl { ( }
          { \tenkz_cd_tree_internal: }
          { \tenkz_cd_tree_leaf: }
      }
      { \PackageError{tenkz}{Malformed~tree:~unexpected~end~of~word}{} }
  }

% a leaf: a wire endpoint on the top line, its label above the tip
\cs_new_protected:Npn \tenkz_cd_tree_leaf:
  {
    \dim_set:Nn \l__tenkzcd_xa_dim
      { \tenkz@r@treestep\tenkz@pitch * \l__tenkzcd_leaf_int }
    \int_incr:N \l__tenkzcd_leaf_int
    \node[label, anchor=south] at
      (\dim_use:N \l__tenkzcd_xa_dim, \tenkz@dim{\tenkz@r@labelclear})
      { $\tenkz@labelsize \tl_use:N \l__tenkzcd_item_tl$ };
    \seq_put_right:Ne \l__tenkzcd_stack_seq
      { { \dim_use:N \l__tenkzcd_xa_dim }{ 0 }{ } }
  }

% ( node node ) [_ charge]
\cs_new_protected:Npn \tenkz_cd_tree_internal:
  {
    \tenkz_cd_tree_node:
    \tenkz_cd_tree_node:
    % the closing parenthesis
    \seq_pop_left:NNTF \l__tenkzcd_items_seq \l__tenkzcd_item_tl
      {
        \str_if_eq:VnF \l__tenkzcd_item_tl { ) }
          { \PackageError{tenkz}
              {Malformed~tree:~expected~a~closing~parenthesis}{} }
      }
      { \PackageError{tenkz}{Malformed~tree:~unexpected~end~of~word}{} }
    % the optional charge (parsed after the recursion, so the local
    % charge register belongs to this vertex).  Peek at the next item: a
    % leading `_` consumes the charge; anything else is the next node and
    % is left in place for the parser to pick up.
    \tl_clear:N \l__tenkzcd_charge_tl
    \seq_get_left:NNT \l__tenkzcd_items_seq \l__tenkzcd_item_tl
      {
        \str_if_eq:VnT \l__tenkzcd_item_tl { _ }
          {
            \seq_pop_left:NN \l__tenkzcd_items_seq \l__tenkzcd_item_tl
            \seq_pop_left:NNF \l__tenkzcd_items_seq \l__tenkzcd_charge_tl
              { \PackageError{tenkz}{Malformed~tree:~dangling~charge}{} }
          }
      }
    % pop the two children and fuse them
    \seq_pop_right:NN \l__tenkzcd_stack_seq \l__tenkzcd_cb_tl
    \seq_pop_right:NN \l__tenkzcd_stack_seq \l__tenkzcd_ca_tl
    \exp_last_unbraced:NV \tenkz_cd_unpack_a:nnn \l__tenkzcd_ca_tl
    \exp_last_unbraced:NV \tenkz_cd_unpack_b:nnn \l__tenkzcd_cb_tl
    \tenkz_cd_tree_fuse:
  }

\cs_new_protected:Npn \tenkz_cd_unpack_a:nnn #1#2#3
  {
    \dim_set:Nn \l__tenkzcd_xa_dim {#1}
    \int_set:Nn \l__tenkzcd_la_int {#2}
    \tl_set:Nn  \l__tenkzcd_cha_tl {#3}
  }
\cs_new_protected:Npn \tenkz_cd_unpack_b:nnn #1#2#3
  {
    \dim_set:Nn \l__tenkzcd_xb_dim {#1}
    \int_set:Nn \l__tenkzcd_lb_int {#2}
    \tl_set:Nn  \l__tenkzcd_chb_tl {#3}
  }

% one child edge: from (x register #1, level #2) up to the shared junction
% at (xp, lp), which the caller has already set
\cs_new_protected:Npn \tenkz_cd_childedge:Nn #1#2
  {
    \draw[bond]
      (\dim_use:N #1, \tenkz_cd_levely:n {#2}) --
      (\dim_use:N \l__tenkzcd_xp_dim, \tenkz_cd_levely:n {\l__tenkzcd_lp_int});
  }

% join the two children on the stack: edges, junction dot, the
% children's charges (their downward edges are now known), push the
% joint carrying this vertex's own charge
\cs_new_protected:Npn \tenkz_cd_tree_fuse:
  {
    \int_set:Nn \l__tenkzcd_lp_int
      { 1 + \int_max:nn { \l__tenkzcd_la_int } { \l__tenkzcd_lb_int } }
    \dim_set:Nn \l__tenkzcd_xp_dim
      { ( \l__tenkzcd_xa_dim + \l__tenkzcd_xb_dim ) / 2 }
    \tenkz_cd_childedge:Nn \l__tenkzcd_xa_dim {\l__tenkzcd_la_int}
    \tenkz_cd_childedge:Nn \l__tenkzcd_xb_dim {\l__tenkzcd_lb_int}
    \node[tree~junction] at
      (\dim_use:N \l__tenkzcd_xp_dim,
       \tenkz_cd_levely:n {\l__tenkzcd_lp_int}) {};
    \tl_if_blank:VF \l__tenkzcd_cha_tl
      { \tenkz_cd_edge_charge:NnN \l__tenkzcd_xa_dim
          {\l__tenkzcd_la_int} \l__tenkzcd_cha_tl }
    \tl_if_blank:VF \l__tenkzcd_chb_tl
      { \tenkz_cd_edge_charge:NnN \l__tenkzcd_xb_dim
          {\l__tenkzcd_lb_int} \l__tenkzcd_chb_tl }
    \seq_put_right:Ne \l__tenkzcd_stack_seq
      { { \dim_use:N \l__tenkzcd_xp_dim }{ \int_use:N \l__tenkzcd_lp_int }
        { \exp_not:V \l__tenkzcd_charge_tl } }
  }

% A charge labels its vertex's downward edge, half a fusion level below
% the vertex, on the side of the edge away from the parent.  Both child
% edges of the vertex lie above it and every other edge lies beyond the
% downward edge on the parent's side, so that quadrant is ink-free by
% construction.  #1 = x register of the vertex, #2 = its level, #3 =
% charge register; the parent's x register is still live.
\cs_new_protected:Npn \tenkz_cd_edge_charge:NnN #1#2#3
  {
    % rightward-descending edge (xp right of the vertex): label the free
    % LEFT side, anchor east, clearance subtracted; leftward or straight:
    % label the RIGHT side, anchor west, clearance added
    \dim_compare:nNnTF { \l__tenkzcd_xp_dim } > { #1 }
      { \tl_set:Nn \l__tenkzcd_anchor_tl {east}
        \dim_set:Nn \l__tenkzcd_off_dim { -\tenkz@dim{\tenkz@r@labelclear} } }
      { \tl_set:Nn \l__tenkzcd_anchor_tl {west}
        \dim_set:Nn \l__tenkzcd_off_dim { \tenkz@dim{\tenkz@r@labelclear} } }
    \node[label, anchor=\l__tenkzcd_anchor_tl] at
      (\dim_eval:n{ #1 + \l__tenkzcd_off_dim },
       \dim_eval:n{ \tenkz_cd_levely:n {#2}
                    - \tenkz@r@treestep\tenkz@pitch / 2 })
      { $\tenkz@labelsize \tl_use:N #3$ };
  }

% the root: a short downward leg; the total charge sits right of the
% leg's midpoint
\cs_new_protected:Npn \tenkz_cd_tree_root:nnn #1#2#3
  {
    \draw[bond]
      (#1, \tenkz_cd_levely:n {#2}) --
      (#1, \dim_eval:n{ \tenkz_cd_levely:n {#2}
                        - \tenkz@r@treeleg\tenkz@pitch });
    \tl_if_blank:nF {#3}
      {
        \node[label, anchor=west] at
          (\dim_eval:n{ #1 + \tenkz@dim{\tenkz@r@labelclear} },
           \dim_eval:n{ \tenkz_cd_levely:n {#2}
                        - \tenkz@r@treeleg\tenkz@pitch / 2 })
          { $\tenkz@labelsize #3$ };
      }
  }

\NewDocumentCommand \tntree { O{} m }
  { \tenkz_cd_tree:nn {#1} {#2} }

\cs_new_protected:Npn \tenkz_cd_tree:nn #1#2
  {
    \group_begin:
    \tl_if_empty:nF {#1} { \pgfqkeys{/tenkz}{#1} }
    % \, and spaces are separators, not syntax; flatten to items
    % (spaces never form items, so only \, needs stripping)
    \tl_set:Nn \l__tenkzcd_word_tl {#2}
    \tl_remove_all:Nn \l__tenkzcd_word_tl { \, }
    \seq_clear:N \l__tenkzcd_items_seq
    \tl_map_inline:Nn \l__tenkzcd_word_tl
      { \seq_put_right:Nn \l__tenkzcd_items_seq {##1} }
    \seq_clear:N \l__tenkzcd_stack_seq
    \int_zero:N \l__tenkzcd_leaf_int
    \tenkz_cd_read_axis:
    \begin{tikzpicture}[ tenkz~every~picture, tenkz~cd~baseline ]
      \tenkz_cd_tree_node:
      \seq_pop_right:NNT \l__tenkzcd_stack_seq \l__tenkzcd_ca_tl
        { \exp_last_unbraced:NV \tenkz_cd_tree_root:nnn \l__tenkzcd_ca_tl }
      % inside the picture group, where the leaf count is still live
      \tenkz@event{tree|picture=\the\tenkz@pictureid|leaves=\int_use:N \l__tenkzcd_leaf_int}
    \end{tikzpicture}
    \group_end:
  }

% ---------- annotation arrows ----------
% \tnarrow[from=i, to=j]{label} records an arrow; the polygon renderer
% draws all recorded arrows once every vertex exists.  The star form
% swaps the label to the other side.  Recording is global because the
% commands sit inside a cell's box.
\NewDocumentCommand \tnarrow { s O{} m }
  {
    \unskip
    \group_begin:
    \tl_clear:N \l__tenkzcd_afrom_tl
    \tl_clear:N \l__tenkzcd_ato_tl
    \tl_if_empty:nF {#2} { \pgfqkeys{/tenkz/arrow}{#2} }
    \seq_gput_right:Ne \g__tenkzcd_arrows_seq
      { { \tl_use:N \l__tenkzcd_afrom_tl }{ \tl_use:N \l__tenkzcd_ato_tl }
        { auto \IfBooleanT {#1} { , swap } }{ \exp_not:n {#3} } }
    \group_end:
    \ignorespaces
  }

\cs_new_protected:Npn \tenkz_cd_draw_arrow:nnnn #1#2#3#4
  {
    % the label's inner sep sets its gap from the shaft; the house 1pt
    % is too tight for sub/superscripted names
    \draw[cd~arrow] (v#1) --
      node[label, inner~sep=2.5pt, #3] { $\tenkz@labelsize #4$ } (v#2);
    \tenkz@event{cdarrow|picture=\the\tenkz@pictureid|from=#1|to=#2}
  }

% ---------- the environment ----------
\NewDocumentEnvironment { tenkzcd } { O{} +b }
  { \tenkz_cd_render:nn {#1} {#2} } { }

\cs_new_protected:Npn \tenkz_cd_render:nn #1#2
  {
    \group_begin:
    \int_zero:N \l__tenkzcd_ngon_int
    \tl_set:Nn \l__tenkzcd_radius_tl {30mm}
    \tl_clear:N \l__tenkzcd_pass_tl
    \tl_if_empty:nF {#1} { \pgfqkeys{/tenkz/cd}{#1} }
    \seq_gclear:N \g__tenkzcd_arrows_seq
    \tenkz@beginpicture{cd}
    \tenkz_cd_read_axis:
    \int_compare:nNnTF { \l__tenkzcd_ngon_int } > {0}
      { \exp_args:NV \tenkz_cd_polygon:nn \l__tenkzcd_pass_tl {#2} }
      { \exp_args:NV \tenkz_cd_grid:nn    \l__tenkzcd_pass_tl {#2} }
    \group_end:
  }

% polygon mode: cell i on the circle at 90 - (i-1) 360/n degrees.  Cells
% render on the compact profile so a default-radius polygon of trees
% leaves room for the arrows; a cell overrides via its own keys.
\cs_new_protected:Npn \tenkz_cd_polygon:nn #1#2
  {
    \dim_set:Nn \l__tenkzcd_r_dim { \l__tenkzcd_radius_tl }
    \seq_set_split:Nnn \l__tenkzcd_cells_seq { & } {#2}
    \begin{tikzpicture}[ tenkz~every~picture, tenkz~cd~baseline, #1 ]
      \int_zero:N \l__tenkzcd_i_int
      \seq_map_inline:Nn \l__tenkzcd_cells_seq
        {
          \int_incr:N \l__tenkzcd_i_int
          % vertex angle: 1 at the top (90 degrees), then clockwise
          \fp_set:Nn \l__tenkzcd_angle_fp
            { 90 - 360 * ( \l__tenkzcd_i_int - 1 ) / \l__tenkzcd_ngon_int }
          \dim_set:Nn \l__tenkzcd_vx_dim
            { \fp_to_decimal:n { cosd( \l__tenkzcd_angle_fp ) } \l__tenkzcd_r_dim }
          \dim_set:Nn \l__tenkzcd_vy_dim
            { \fp_to_decimal:n { sind( \l__tenkzcd_angle_fp ) } \l__tenkzcd_r_dim }
          \node[inner~sep=3pt] (v\int_use:N \l__tenkzcd_i_int) at
            (\dim_use:N \l__tenkzcd_vx_dim, \dim_use:N \l__tenkzcd_vy_dim)
            { \pgfkeys{/tenkz/compact} ##1 };
          \tenkz@event{cdcell|picture=\the\tenkz@pictureid|index=\int_use:N \l__tenkzcd_i_int}
        }
      \seq_map_inline:Nn \g__tenkzcd_arrows_seq
        { \tenkz_cd_draw_arrow:nnnn ##1 }
    \end{tikzpicture}
  }

% Grid mode is a thin wrapper around tikz-cd: cells hold tenkz objects
% or mathematics, arrows use tikz-cd's own \arrow syntax.  The house
% annotation line width is applied to every arrow; arrow tips stay
% tikz-cd's defaults (restyling the computed tip is not worth the
% coupling in Phase 0).
%
% The body reaches us already tokenized, so pgfmatrix's catcode trick
% for & never fires; every & is rewritten to the cell separator it
% would have become.  (Consequence: a raw & nested inside a cell needs
% its own guard, e.g. a group with ampersand replacement.)
\cs_new_protected:Npn \tenkz_cd_grid:nn #1#2
  {
    \tl_set:Nn \l__tenkzcd_body_tl {#2}
    \tl_replace_all:Nnn \l__tenkzcd_body_tl { & } { \pgfmatrixnextcell }
    \tikzcdset{ every~arrow/.append~style =
                  { line~width = \tenkz@annwidth } }
    \exp_args:NnV \tenkz_cd_grid_aux:nn {#1} \l__tenkzcd_body_tl
  }
\cs_new_protected:Npn \tenkz_cd_grid_aux:nn #1#2
  {
    % same baseline convention as the `tenkz cd baseline` style, inlined
    % because tikzcd routes its options through its own key family
    \begin{tikzcd}
      [ baseline={([yshift=-\dim_use:N \l__tenkzcd_axis_dim]
                   current~bounding~box.center)}, #1 ]
      #2
    \end{tikzcd}
  }

\ExplSyntaxOff
\endinput
